%==============================================================================
%  MATH 231 -- Linear Algebra I (Queens College, Fall 2026)
%  SETUP GUIDE 2 -- Python, pip, and an IDE
%
%  AUDIENCE: students, some of whom have never installed a programming language.
%  Companion to Setup Guide 1 (git and GitHub).  Same conventions: numbered
%  steps, per-platform instructions, a Checkpoint after every step, and a
%  troubleshooting section keyed to real error messages.
%
%  The package list is deliberately short.  numpy is the only one installed
%  here, as practice; scipy and matplotlib arrive when the course needs them.
%
%  Compile with pdflatex (standard TeX Live packages only).
%==============================================================================
\documentclass[11pt]{article}

\usepackage[margin=1in]{geometry}
\usepackage{amsmath,amssymb}
\usepackage[dvipsnames]{xcolor}
\usepackage{enumitem}
\usepackage{booktabs}
\usepackage{array}
\usepackage[hidelinks]{hyperref}
\usepackage{titlesec}
\usepackage{needspace}
\usepackage{listings}

%------------------------------------------------------------------ style ----
\titleformat{\section}{\large\bfseries\sffamily\color{RoyalBlue}}
                      {\thesection.}{0.6em}{}
\titleformat{\subsection}{\normalsize\bfseries\sffamily}{\thesubsection.}{0.5em}{}
\titlespacing*{\section}{0pt}{16pt}{7pt}

\lstdefinestyle{term}{
  basicstyle=\ttfamily\small,
  backgroundcolor=\color{gray!10},
  frame=single, rulecolor=\color{gray!45},
  framesep=5pt, xleftmargin=8pt, xrightmargin=8pt,
  breaklines=true, columns=fullflexible, keepspaces=true,
  showstringspaces=false, aboveskip=8pt, belowskip=8pt
}
\lstdefinestyle{py}{
  style=term, language=Python,
  keywordstyle=\color{RoyalBlue}\bfseries,
  commentstyle=\color{gray!80}\itshape,
  stringstyle=\color{ForestGreen},
  emph={np,numpy}, emphstyle=\color{BrickRed}
}
\lstset{style=term}

\newcommand{\plat}[1]{\smallskip\noindent
  \textcolor{BrickRed}{\sffamily\bfseries #1}\par\nopagebreak\smallskip}
\newcommand{\checkpoint}[1]{\par\smallskip\noindent
  {\small\sffamily\textcolor{ForestGreen}{\textbf{Checkpoint.}}\ #1}\par\smallskip}
\newcommand{\heads}[1]{\par\smallskip\noindent
  {\small\sffamily\textcolor{BrickRed}{\textbf{Watch out.}}\ #1}\par\smallskip}

\title{\vspace{-1.2cm}\textbf{\sffamily MATH 231 -- Linear Algebra I}\\[4pt]
       {\large\sffamily Setup Guide 2 \textbullet\ Python, pip, and an IDE}\\[2pt]
       {\normalsize\sffamily Queens College, CUNY \textbullet\ Fall 2026
        \textbullet\ Instructor: Kevin Bedoya}}
\author{}
\date{}

\begin{document}
\maketitle
\vspace{-1.4cm}

\noindent\textbf{What this guide does.} By the end you will have Python
installed, the package installer \texttt{pip} working, a code editor set up, the
\texttt{numpy} package installed, and a short program that runs and prints two
vectors. That is the complete toolkit for every piece of programming in this
course.

\medskip
\noindent\textbf{No experience assumed.} If you have never written a line of
code, you are in the right place --- Python is genuinely one of the friendlier
languages to start with, and programming will never be assessed on an exam in
this course. Follow the steps in order and check the
\textcolor{ForestGreen}{\sffamily\bfseries Checkpoint} after each. Section \S8
is troubleshooting.

\medskip
\noindent\textbf{Everything here is free.} No licences, no trials, no credit
card.

%==============================================================================
\section{What you are installing, and why}
%==============================================================================
Three separate pieces. They are often confused, so it is worth naming them.

\begin{center}
\renewcommand{\arraystretch}{1.5}
\begin{tabular}{@{}l p{10.2cm}@{}}
\toprule
\textbf{Piece} & \textbf{What it is} \\
\midrule
\textbf{Python} & The language itself, plus the program that runs your code.
                  Without it, a \texttt{.py} file is just text. \\
\textbf{pip}    & Python's package installer. It fetches libraries other people
                  have written --- like \texttt{numpy} --- and puts them where
                  Python can find them. It comes bundled with Python. \\
\textbf{An IDE} & Where you write the code. Technically optional, but it gives
                  you syntax colouring, error highlighting, and a Run button. \\
\bottomrule
\end{tabular}
\end{center}

\noindent Install them in that order: Python first, since pip arrives with it,
and the IDE last, since it needs to find the Python you installed.

%==============================================================================
\section{Install Python}
%==============================================================================
Find your operating system.

\plat{Windows}
\begin{enumerate}[leftmargin=2.0em,itemsep=4pt]
  \item Go to \url{https://www.python.org/downloads/} and click the big
        \textbf{Download Python 3.x} button.
  \item Run the installer. On the very first screen, \textbf{tick the box at the
        bottom that says ``Add python.exe to PATH''}.
  \item Then click \textbf{Install Now} and let it finish.
\end{enumerate}

\heads{That checkbox is the single most important step in this guide. It is off
by default, it is easy to miss, and skipping it is the cause of roughly every
``Python is not recognised'' message. If you are not sure whether you ticked it,
run the installer again and choose \emph{Modify} --- reinstalling does no
harm.}

\plat{macOS}
Go to \url{https://www.python.org/downloads/} and run the macOS installer,
accepting the defaults. Alternatively, if you already use Homebrew:

\begin{lstlisting}
brew install python
\end{lstlisting}

\heads{macOS ships with an old Python that the system uses for its own purposes.
Leave it alone. On a Mac, always type \texttt{python3} and \texttt{pip3} rather
than \texttt{python} and \texttt{pip}.}

\plat{Linux}
Python 3 is usually already installed. Check first, and install only if it is
missing:

\begin{lstlisting}
sudo apt update && sudo apt install python3 python3-pip   # Ubuntu, Debian, Mint
sudo dnf install python3 python3-pip                      # Fedora
\end{lstlisting}

\subsection*{Check that it worked}
Open a \textbf{new} terminal window --- new, because an already-open one will not
know about a program installed a minute ago --- and run:

\begin{lstlisting}
python --version        # Windows
python3 --version       # macOS and Linux
\end{lstlisting}

\checkpoint{You see something like \texttt{Python 3.12.4}. Any version starting
with \texttt{3.} is fine. If it says \texttt{2.7}, see \S8.4.}

%==============================================================================
\section{Check pip}
%==============================================================================
There is nothing to install --- pip ships with every modern Python. You only
need to confirm it is there:

\begin{lstlisting}
python -m pip --version        # Windows
python3 -m pip --version       # macOS and Linux
\end{lstlisting}

\noindent The \texttt{-m} means \emph{run the module named pip using this
Python}. It is slightly longer to type than plain \texttt{pip}, and worth the
habit: it guarantees that packages get installed into the same Python that runs
your code. Most beginner package problems come from those two being different.

\checkpoint{You see a line like \texttt{pip 24.0 from ... (python 3.12)}.}

\medskip
\noindent Optionally, bring pip up to date:

\begin{lstlisting}
python -m pip install --upgrade pip
\end{lstlisting}

%==============================================================================
\section{Install an IDE}
%==============================================================================
An \textbf{IDE} (integrated development environment) is a text editor built for
code. Pick \emph{one} of the two below --- both are free, both are widely used
in industry, and either is completely fine for this course.

\begin{center}
\renewcommand{\arraystretch}{1.5}
\begin{tabular}{@{}l p{9.6cm}@{}}
\toprule
\textbf{Option} & \textbf{Choose it if} \\
\midrule
\textbf{VS Code} & You want something light and general-purpose that also
                   handles other languages, or you already use it. \\
\textbf{PyCharm} & You want everything preconfigured for Python specifically,
                   with less setup and more built in. \\
\bottomrule
\end{tabular}
\end{center}

\subsection{VS Code}
\begin{enumerate}[leftmargin=2.0em,itemsep=4pt]
  \item Download from \url{https://code.visualstudio.com} and install it.
  \item Open it and click the \textbf{Extensions} icon in the left sidebar (four
        small squares).
  \item Search for \textbf{Python} and install the one published by
        \textbf{Microsoft}. This is what makes VS Code understand Python rather
        than treating your file as plain text.
  \item Press \texttt{Ctrl}+\texttt{Shift}+\texttt{P}
        (\texttt{Cmd}+\texttt{Shift}+\texttt{P} on a Mac), type
        \textbf{Python: Select Interpreter}, and choose the Python 3 you
        installed in \S2.
\end{enumerate}

\subsection{PyCharm}
\begin{enumerate}[leftmargin=2.0em,itemsep=4pt]
  \item Download from \url{https://www.jetbrains.com/pycharm/download/}.
  \item Choose the \textbf{Community Edition} --- it is free and has everything
        this course needs. You do not need the paid Professional edition.
  \item Install and open it. When you create a new project, PyCharm will find
        your Python installation on its own.
\end{enumerate}

\checkpoint{You can open your editor, create a file called \texttt{test.py}, and
see the code coloured rather than uniformly black.}

%==============================================================================
\section{Install your first package: numpy}
%==============================================================================
A \textbf{package} is code somebody else wrote that you can use in your own
programs. \texttt{numpy} --- ``Numerical Python'' --- is the standard package
for numerical work, and it is the one we will lean on most this semester, since
it is built around exactly the object this course is about: the vector.

\medskip
\noindent Install it:

\begin{lstlisting}
python -m pip install numpy        # Windows
python3 -m pip install numpy       # macOS and Linux
\end{lstlisting}

\noindent It will print a few lines of download progress and finish with a line
beginning \texttt{Successfully installed numpy}.

\checkpoint{Run the line below. It should print a version number and no errors.}

\begin{lstlisting}
python -c "import numpy; print(numpy.__version__)"
\end{lstlisting}

\medskip
\noindent \textbf{Later packages.} We will add others as the course needs them
--- \texttt{scipy} and \texttt{matplotlib} are the likely next two --- and each
is installed the same way, by swapping the name in the command above. There is
nothing to install ahead of time; I will say when.

%==============================================================================
\section{Your first program}
%==============================================================================
Time to check that all three pieces work together.

\begin{enumerate}[leftmargin=2.0em,itemsep=4pt]
  \item In your editor, create a new file named \texttt{main.py}.
  \item Type the following into it. Three lines of work, plus the import that
        makes \texttt{numpy} available.
\end{enumerate}

\begin{lstlisting}[style=py]
import numpy as np
u = np.array([1, 2, 3])
v = np.array([4, 5, 6])
print(u, v)
\end{lstlisting}

\noindent Line by line:

\begin{itemize}[leftmargin=1.4em,itemsep=4pt,topsep=4pt]
  \item \texttt{import numpy as np} loads the package and gives it the short
        nickname \texttt{np}. Essentially all Python code in the world writes it
        this way, so it is worth copying the convention.
  \item \texttt{np.array([1, 2, 3])} builds a numpy array --- which for our
        purposes is a vector --- from an ordinary list of numbers. We store one
        in \texttt{u} and another in \texttt{v}.
  \item \texttt{print(u, v)} writes both to the terminal.
\end{itemize}

\noindent Save the file, then run it. Either press the \textbf{Run} button in
your IDE, or from a terminal in the same folder as the file:

\begin{lstlisting}
python main.py         # Windows
python3 main.py        # macOS and Linux
\end{lstlisting}

\checkpoint{The terminal prints exactly this:}

\begin{lstlisting}
[1 2 3] [4 5 6]
\end{lstlisting}

\noindent Note that numpy separates components with spaces rather than commas,
which is how you can tell a numpy array from an ordinary Python list at a
glance.

\medskip
\noindent If you see that output, your setup is complete: Python runs, pip
installed a package, the package imported, and your editor can execute a file.
Nothing further is needed until the first assignment.

%==============================================================================
\section{Where your code goes}
%==============================================================================
Assignments are submitted through the private GitHub repository set up in
\textbf{Setup Guide 1}. Save your \texttt{.py} files into the matching folder
--- \texttt{homework/}, \texttt{projects/}, or \texttt{extra-credit/} --- inside
your \texttt{QC\_MATH231} folder, then stage, commit, and push:

\begin{lstlisting}
git add .
git commit -m "Finish homework 1"
git push
\end{lstlisting}

\noindent Work that has not been pushed is not visible to me and cannot be
graded. If any of that is unfamiliar, it is covered in Setup Guide 1, \S7.

%==============================================================================
\section{Troubleshooting}
%==============================================================================

\subsection{``python is not recognized'' / command not found}
\begin{lstlisting}
'python' is not recognized as an internal or external command
\end{lstlisting}
Almost always the \emph{Add python.exe to PATH} checkbox from \S2. Two fixes, in
order of effort:

\begin{enumerate}[leftmargin=2.0em,itemsep=3pt]
  \item \textbf{Open a new terminal.} An open terminal does not pick up a
        program installed after it started.
  \item \textbf{Try the launcher.} On Windows, \texttt{py} often works even when
        \texttt{python} does not:
\begin{lstlisting}
py --version
\end{lstlisting}
        If that works, use \texttt{py} in place of \texttt{python} everywhere,
        including \texttt{py -m pip install numpy}.
  \item \textbf{Reinstall.} Run the installer again, choose \textbf{Modify}, and
        make sure the PATH option is ticked.
\end{enumerate}

\subsection{Typing \texttt{python} opens the Microsoft Store}
A Windows default that intercepts the command when no real Python is on the
PATH. Either install Python properly from \url{https://python.org} with the PATH
box ticked, or turn the stub off: \textbf{Settings} $\to$ \textbf{Apps} $\to$
\textbf{Advanced app settings} $\to$ \textbf{App execution aliases}, then switch
off both \texttt{python.exe} and \texttt{python3.exe}.

\subsection{``pip is not recognized''}
Use the longer form, which works whenever Python itself works:

\begin{lstlisting}
python -m pip install numpy
\end{lstlisting}

\subsection{\texttt{python --version} says 2.7}
You are on a Mac or an older Linux system where \texttt{python} means the legacy
version. Use \texttt{python3} and \texttt{pip3} instead --- everywhere, for the
whole semester. Python 2 will not run this course's code.

\subsection{``ModuleNotFoundError: No module named numpy''}
\begin{lstlisting}
ModuleNotFoundError: No module named 'numpy'
\end{lstlisting}
numpy went into a different Python than the one running your file. This is the
most common cause of confusion when more than one Python is installed. Fix it by
installing through the same interpreter that runs your code:

\begin{lstlisting}
python -m pip install numpy
\end{lstlisting}

\noindent If it still fails inside your IDE but works in the terminal, the IDE is
pointed at the wrong interpreter. In VS Code, press
\texttt{Ctrl}+\texttt{Shift}+\texttt{P} $\to$ \textbf{Python: Select
Interpreter} and choose the same Python 3 you have been using. In PyCharm:
\textbf{Settings} $\to$ \textbf{Project} $\to$ \textbf{Python Interpreter}.

\subsection{Permission denied while installing}
Do not reach for \texttt{sudo}. Install into your own user account instead:

\begin{lstlisting}
python -m pip install --user numpy
\end{lstlisting}

\subsection{``SyntaxError'' on a line that looks fine}
Check the line \emph{above} the one reported --- an unclosed bracket or quote
makes Python complain about the next line down. Also check that you have not
mixed tabs and spaces for indentation; most editors have a setting to insert
spaces for tabs.

\subsection{The program runs but prints nothing}
You probably saved the file without the \texttt{print} line, or ran a different
file than the one you edited. Confirm the filename in the terminal output
matches the file you are looking at, and that you saved before running.

\subsection{Still stuck}
Email me at \url{Kevin.Bedoya32@stu-mail.qc.cuny.edu} or ask on the class
Discord. Include:

\begin{itemize}[leftmargin=1.4em,itemsep=2pt,topsep=3pt]
  \item your operating system,
  \item the exact command you ran, and
  \item the full error message, copied as text or as a screenshot.
\end{itemize}

\noindent Installation problems are ordinary and everyone hits at least one.
Sort yours out early --- well before the first assignment is due --- so that
when the deadline arrives the only thing you are thinking about is the linear
algebra.

\end{document}
